% 《物理实验报告模板.tex》的共同版式；所有实验模板均载入本文件。
\documentclass[UTF8,a4paper,10pt,fontset=fandol]{ctexart}
\usepackage[left=3.175cm,right=3.175cm,top=2.54cm,bottom=2.54cm]{geometry}
\usepackage{graphicx,tikz,pgfplots,needspace,adjustbox,booktabs,fancyhdr,lastpage,amsmath,amssymb,array,tabularx,longtable}
\pgfplotsset{compat=1.18}
\setlength{\parindent}{0pt}
\setlength{\parskip}{2pt}
\setlength{\headheight}{18pt}
\linespread{1}
\newcommand{\SetField}[2]{\expandafter\def\csname reportfield@#1\endcsname{#2}}
\makeatletter
\newcommand{\Field}[2][1.35cm]{%
  \ifcsname reportfield@#2\endcsname
    \expandafter\ifx\csname reportfield@#2\endcsname\@empty
      \underline{\makebox[#1]{\strut}}%
    \else\csname reportfield@#2\endcsname\fi
  \else\underline{\makebox[#1]{\strut}}\fi}
\makeatother
\newcommand{\MathField}[2][1.35cm]{%
  \ifcsname reportfield@math_#2\endcsname
    \csname reportfield@math_#2\endcsname
  \else\text{\Field[#1]{#2}}\fi}
\newcommand{\coverfield}[2]{%
  \noindent{\heiti\fontsize{18pt}{48pt}\selectfont
  \makebox[3.3cm][l]{#1}\underline{\makebox[8.8cm][c]{\Field[0pt]{#2}}}\par}}
\newcommand{\bindingline}[1][2.85]{%
  \begin{tikzpicture}[remember picture,overlay]
    \foreach \i in {0,...,38}{%
      \pgfmathsetmacro{\bindingy}{-#1-0.52*\i}%
      \ifnum\i=13\relax
        \node[inner sep=0pt] at ([xshift=1.05cm,yshift=\bindingy cm]current page.north west) {\songti\fontsize{10.5pt}{10.5pt}\selectfont 装};
      \else\ifnum\i=19\relax
        \node[inner sep=0pt] at ([xshift=1.05cm,yshift=\bindingy cm]current page.north west) {\songti\fontsize{10.5pt}{10.5pt}\selectfont 订};
      \else\ifnum\i=25\relax
        \node[inner sep=0pt] at ([xshift=1.05cm,yshift=\bindingy cm]current page.north west) {\songti\fontsize{10.5pt}{10.5pt}\selectfont 线};
      \else
        \pgfmathsetmacro{\bindingtop}{\bindingy+0.20}%
        \pgfmathsetmacro{\bindingbottom}{\bindingy-0.20}%
        \draw[line width=0.35pt,dash pattern=on 1.2pt off 1.8pt]
          ([xshift=1.05cm,yshift=\bindingtop cm]current page.north west)
          -- ([xshift=1.05cm,yshift=\bindingbottom cm]current page.north west);
      \fi\fi\fi}
  \end{tikzpicture}%
}
% 装订线印在每一页正文和原始数据页，续页也沿用母模板样式。
\newif\ifreportbinding
\AddToHook{shipout/foreground}{\ifreportbinding\bindingline\fi}
\fancypagestyle{report}{\fancyhf{}\fancyfoot[C]{第\thepage 页\quad 共\pageref{LastPage}页}\renewcommand{\headrulewidth}{0.4pt}}
\fancypagestyle{record}{\fancyhf{}\fancyhead[C]{原始数据记录单}\fancyfoot[C]{第\thepage 页\quad 共\pageref{LastPage}页}\renewcommand{\headrulewidth}{0.4pt}}
\newcommand{\ReportCover}[1]{%
  \thispagestyle{empty}\vspace*{1.0cm}
  \begin{center}
    \includegraphics[width=8.184cm]{tongji_logo.jpeg}\par
    {\bfseries\fontsize{18pt}{24pt}\selectfont TONGJI UNIVERSITY\par}
    \vspace{1.5cm}{\heiti\fontsize{36pt}{44pt}\selectfont 物理实验报告\par}
  \end{center}
  \vspace{1.5cm}
  \coverfield{实验名称}{experiment}
  \coverfield{学生姓名}{student}
  \coverfield{学\hspace{1em}号}{student_id}
  \coverfield{组\hspace{1em}号}{group}
  \coverfield{实验日期}{date}
  \clearpage\reportbindingtrue\pagestyle{report}
  \begin{center}{\heiti\fontsize{14pt}{40pt}\selectfont #1\par}\end{center}}
\newcommand{\ReportSection}[1]{\Needspace{4\baselineskip}\par\vspace{0.35\baselineskip}{\heiti #1}\par}
\newcommand{\ReportSubsection}[1]{\Needspace{3\baselineskip}\par\vspace{0.2\baselineskip}{\bfseries #1}\par}
\newif\ifreportimages
\newcommand{\ReportRecord}{%
  \clearpage\thispagestyle{record}
  \vspace*{0.6cm}%
  \ifreportimages\else（请将教师签字的实验数据记录单完整拍照粘贴至此处。）\fi
}

% ===== 数据输入区开始 =====
\SetField{experiment}{阴极射线示波器}
% 后端以 \SetField{字段名}{值} 填入数据；缺省值保持空白。
% 母模板封面字段：\Field{experiment} \Field{student} \Field{student_id} \Field{group} \Field{date}
% ===== 数据输入区结束 =====
\begin{document}\fontsize{10.5pt}{20pt}\selectfont
\ReportCover{阴极射线示波器}
\ReportSection{实验目的}
1. 了解阴极射线示波器的结构、电偏转、扫描和触发同步原理。\par
2. 掌握示波器校准及峰峰值、周期、频率的测量方法。\par
3. 观察不同波形，利用李萨如图形测量正弦信号频率。
\ReportSection{实验仪器}
ST-16 型单踪示波器、SG1001 型简易信号发生器、UTG7025B 型函数信号发生器及同轴电缆。
\ReportSection{实验原理}
\ReportSubsection{电子束偏转与波形显示}
电子枪发射的电子经加速、聚焦后轰击荧光屏。两组互相垂直的偏转板控制亮点的横、纵位移。
在小偏转近似下，位移与偏转电压成正比：$y=D_yU_y$、$x=D_xU_x$。
Y 轴输入待测信号，X 轴输入与时间成正比的锯齿扫描电压，就把电压随时间的变化展开成波形。
触发电路让每次扫描从待测信号的同一相位开始，从而稳定显示周期性波形。
\begin{center}
\begin{tikzpicture}[every node/.style={font=\small},>=stealth]
\node[draw,minimum width=2.2cm,minimum height=.8cm] (v) at (0,1.4) {垂直放大};
\node[draw,minimum width=2.2cm,minimum height=.8cm] (tr) at (0,0) {触发同步};
\node[draw,minimum width=2.2cm,minimum height=.8cm] (sw) at (3,0) {扫描发生器};
\node[draw,minimum width=2.2cm,minimum height=.8cm] (h) at (6,0) {水平放大};
\node[draw,minimum width=2.2cm,minimum height=.8cm] (crt) at (6,1.4) {示波管};
\draw[->] (-2.4,1.4) node[left]{Y 输入} -- (v);
\draw[->] (v) -- (crt); \draw[->] (v) -- (tr);
\draw[->] (tr) -- (sw); \draw[->] (sw) -- (h); \draw[->] (h) -- (crt);
\end{tikzpicture}\par {\small 示波器扫描与触发的简化结构。}
\end{center}
\ReportSubsection{电压、周期与频率}
峰峰值 $U_{\rm pp}$ 由纵向跨度 $N_U$ 与电压档位 $S_y$ 相乘得到；周期由横向跨度 $N_T$ 与时基档位 $S_x$ 相乘得到：
\[U_{\rm pp}=K S_yN_U,\qquad T=S_xN_T,\qquad f=\frac1T.\]
$K$ 为探头衰减倍数，同轴直连时 $K=1$；$S_y$ 用 V/div，$S_x$ 用 ms/div，故 $T$ 用 ms，波形频率 $f_v=1/T$ 用 kHz；李萨如法频率用 kHz。
峰峰值并非峰值或有效值；仅对正弦信号有 $U_{\rm rms}=U_{\rm pp}/(2\sqrt2)$。

\clearpage
\ReportSection{实验内容与操作}
\ReportSubsection{预备调节与校准}
接通电源并预热，调节适中辉度、聚焦及 X/Y 位移，避免亮点长时间停留在同一位置。
Y 输入使用合适耦合；观察交流波形可选 AC，测直流必须选 DC。
校准时使用 X1 探头连接内置标准方波（$0.5\,\mathrm{V_{pp}}$、$1\,\mathrm{kHz}$）。
将垂直档设为 $0.1\,\mathrm{V/div}$、时基设为 $0.1\,\mathrm{ms/div}$，预期峰峰跨度为 5 div、周期跨度为 10 div。
必要时调节灵敏度微调，校准后测量期间保持不变。\par
\ReportSubsection{四种信号波形测量}
将 SG1001 输出接示波器 Y 输入，按原始记录表顺序选择方波、正弦波、三角波、半波。
调节电压档位、时基及 LEVEL，使波形稳定且跨度尽量大，记录峰峰跨度和一个完整周期的跨度。
多周期测量时，应先把总跨度除以周期数，再填入单周期跨度。
\ReportSubsection{李萨如图形测频}
把已知频率 $f_x$ 的正弦信号接 X 输入，待测正弦信号接 Y 输入，切换到外接 X 或 X--Y 模式。
保持 Y 信号频率不变，改变 X 信号频率，取得四组稳定的简单闭合图形。
分别计数图形与假想水平线、垂直线的切点数，记为 $n_x$、$n_y$，不要把普通交点或曲线自交点当作切点。
记录图形草图或拍照，与切点计数一同保存在原始记录中。
\ReportSubsection{注意事项}
避免频繁开关机；测量前校准，档位切换轻缓；示波器与信号源公共地连接一致。
大幅度信号先衰减，测量前确认探头倍数；示波器内有高压，禁止自行打开外壳。

\clearpage
\ReportSection{数据记录与处理：波形测量}
探头倍数 $K=\Field{probe_factor}$；同轴直连取 1。本报告四组数据使用同一探头倍数。
表格顺序与原始记录一致：方波、正弦波、三角波、半波；半波的周期取相邻重复结构的间隔。
\ReportSubsection{电压峰峰值}
\[U_{\rm pp}=K S_yN_U.\]
\begin{center}\small\renewcommand{\arraystretch}{1.35}
\begin{adjustbox}{max width=\textwidth}
\begin{tabular}{|c|c|c|c|c|}\hline
序号 & 波形 & $S_y$/(V/div) & $N_U$/div & $U_{\rm pp}$/V\\\hline
1 & 方波 & \Field[10mm]{voltage_scale_01} & \Field[10mm]{voltage_divisions_01} & \Field[10mm]{voltage_pp_01}\\\hline
2 & 正弦波 & \Field[10mm]{voltage_scale_02} & \Field[10mm]{voltage_divisions_02} & \Field[10mm]{voltage_pp_02}\\\hline
3 & 三角波 & \Field[10mm]{voltage_scale_03} & \Field[10mm]{voltage_divisions_03} & \Field[10mm]{voltage_pp_03}\\\hline
4 & 半波 & \Field[10mm]{voltage_scale_04} & \Field[10mm]{voltage_divisions_04} & \Field[10mm]{voltage_pp_04}\\\hline
\end{tabular}
\end{adjustbox}
\end{center}
\ReportSubsection{周期与频率}
\[T=S_xN_T,\qquad f_v\,(\mathrm{kHz})=\frac{1}{T\,(\mathrm{ms})}.\]
\begin{center}\small\renewcommand{\arraystretch}{1.35}
\begin{adjustbox}{max width=\textwidth}
\begin{tabular}{|c|c|c|c|c|c|}\hline
序号 & 波形 & $S_x$/(ms/div) & $N_T$/div & $T$/ms & $f_v$/kHz\\\hline
1 & 方波 & \Field[10mm]{time_scale_01} & \Field[10mm]{period_divisions_01} & \Field[10mm]{period_ms_01} & \Field[10mm]{frequency_khz_01}\\\hline
2 & 正弦波 & \Field[10mm]{time_scale_02} & \Field[10mm]{period_divisions_02} & \Field[10mm]{period_ms_02} & \Field[10mm]{frequency_khz_02}\\\hline
3 & 三角波 & \Field[10mm]{time_scale_03} & \Field[10mm]{period_divisions_03} & \Field[10mm]{period_ms_03} & \Field[10mm]{frequency_khz_03}\\\hline
4 & 半波 & \Field[10mm]{time_scale_04} & \Field[10mm]{period_divisions_04} & \Field[10mm]{period_ms_04} & \Field[10mm]{frequency_khz_04}\\\hline
\end{tabular}
\end{adjustbox}
\end{center}
\ReportSubsection{读数与误差说明}
读取峰、谷时选取稳定轮廓，读取周期时选取相邻的同相位点。方波和脉冲的边沿有一定宽度，应保持同一读数标准。
时基和垂直微调若离开校准位置，档位标称值不能直接用于测量。
本实验未输入仪器误差限，电压与周期表仅给出由读数得到的结果；不能把显示小数位数当作仪器准确度。

\clearpage
\ReportSection{数据记录与处理：李萨如图形}
频率比由切点数确定：
\[\frac{f_y}{f_x}=\frac{n_x}{n_y},\qquad f_y=\frac{n_x}{n_y}f_x.\]
\begin{center}
\begin{tikzpicture}
\foreach \index/\offset in {01/0,02/3.2,03/6.4,04/9.6} {
\begin{scope}[xshift=\offset cm]
\draw[gray] (-1.25,-1) rectangle (1.25,1);
\ifcsname reportfield@nx_\index\endcsname
\ifcsname reportfield@ny_\index\endcsname
\pgfmathsetmacro{\p}{\csname reportfield@nx_\index\endcsname}
\pgfmathsetmacro{\q}{\csname reportfield@ny_\index\endcsname}
\pgfmathtruncatemacro{\plotSamples}{48*max(\p,\q)+1}
\draw[blue,thin,samples=\plotSamples,domain=0:360,variable=\t]
plot ({1.1*sin(\q*\t+30)},{.85*sin(\p*\t)});
\fi\fi
\node[font=\small,align=center] at (0,-1.45) {第\index 组\\$f_y:f_x=\Field{frequency_ratio_\index}$};
\end{scope}}
\end{tikzpicture}\par
{\small 按各组输入的切点数绘制频率比图；振幅归一化，相位取示意值，并非实测波形。}
\end{center}
\begin{center}\small\renewcommand{\arraystretch}{1.35}
\begin{adjustbox}{max width=\textwidth}
\begin{tabular}{|c|c|c|c|c|c|}\hline
组次 & $f_x$/kHz & $n_y$ & $n_x$ & $f_y:f_x$ & $f_y$/kHz\\\hline
1 & \Field[10mm]{reference_frequency_01} & \Field[10mm]{ny_01} & \Field[10mm]{nx_01} & \Field[10mm]{frequency_ratio_01} & \Field[10mm]{lissajous_frequency_01}\\\hline
2 & \Field[10mm]{reference_frequency_02} & \Field[10mm]{ny_02} & \Field[10mm]{nx_02} & \Field[10mm]{frequency_ratio_02} & \Field[10mm]{lissajous_frequency_02}\\\hline
3 & \Field[10mm]{reference_frequency_03} & \Field[10mm]{ny_03} & \Field[10mm]{nx_03} & \Field[10mm]{frequency_ratio_03} & \Field[10mm]{lissajous_frequency_03}\\\hline
4 & \Field[10mm]{reference_frequency_04} & \Field[10mm]{ny_04} & \Field[10mm]{nx_04} & \Field[10mm]{frequency_ratio_04} & \Field[10mm]{lissajous_frequency_04}\\\hline
\end{tabular}
\end{adjustbox}
\end{center}
\[\bar f_y=\frac14\sum_{i=1}^4f_{y,i},\quad
u_A(\bar f_y)=\sqrt{\frac{\sum_{i=1}^4(f_{y,i}-\bar f_y)^2}{4(4-1)}}.\]
\noindent 李萨如法：$\bar f_y=(\MathField{lissajous_mean_khz}\pm\MathField{lissajous_u_khz})\,\mathrm{kHz}$。
这里只估计重复测频的 A 类标准不确定度，未计入参考频率准确度。
\par 正弦波周期法：$f_T=\MathField{sine_frequency_khz}\,\mathrm{kHz}$；
两法相对差 $|\bar f_y-f_T|/f_T=\Field{method_difference_percent}\%$。
平均与比较要求四组记录和第 2 行正弦波测量使用同一个、频率未改变的待测信号。
\ReportSection{讨论与分析（思考题）}
1. 只见竖直亮线，通常是 X 扫描未开启或未接入合适扫描信号。应检查时基、X 输入选择及水平放大，调节触发源、电平和档位，使波形展开并稳定。\par
2. 测直流时用 DC 耦合，先确定零电压基线，再接入待测电压；电压为基线位移乘垂直档位及探头倍数，正负由偏转方向确定。\par
3. 普通触发扫描难以持续稳定显示非周期信号，因为重复扫描不能保持同一相位；瞬态捕获需要单次扫描或存储示波器。\par
4. 李萨如图形转动或变化表明频率比未稳定或相位缓慢漂移。微调信号频率、检查连接并保持信号源稳定，可得到便于计数的闭合图形。
\ReportSection{实验结论}
通过校准和档位读数得到四种波形的峰峰值、周期及频率。
李萨如法平均频率为 $\MathField{lissajous_mean_khz}\,\mathrm{kHz}$，与周期法的相对差为 \Field{method_difference_percent}\%。
评价两法一致性时应结合切点计数、读数分辨率和参考频率准确度；上述相对差本身不是准确度判定标准。

\ReportRecord
% ===== 原始数据图片插入区 =====
\end{document}
